% s-088-p088-m003; original PDF 88, 89
\Aphorism{s-088-p088-m003}{89}
In alten Gesetzgebungen hiess es, dem Verleumder solle die Zunge ausgeschnitten, dem Diebe die Hand abgehauen werden; jeder Sünder solle mit dem Gliede büssen, mit dem er gesündigt. Das muthet uns heute naiv und oberflächlich-realistisch an; wir wissen, hinter den Extremitäten sitzen Muskeln, hinter diesen wieder Nerven — die Lästerzunge und die Diebeshand ist nur das letzte, in die Sichtbarkeit hinaustretende Glied eines langen Complexes von Ursache und Wirkung. Im Grunde aber, hegen wir nicht noch denselben Aberglauben, indem wir überhaupt den Verbrecher strafen? Sollte nicht auch im socialen Organismus hinter dem Verbrecher ein Agens sitzen, eine uns noch unbekannte Art von Wille, deren Hervortreten ins Augenscheinliche eben Verbrechen heisst? Und unser Strafen wäre demnach Rache an einem Executivorgan niedersten Ranges, statt am eigentlich Schuldigen, — unsere Themis wäre wirklich blind?
\par

